% Section 13.7, from lectures/L13/README.md.

\section{Summary}
\label{c13:sec:review}

\needspace{6\baselineskip}
\subsection*{What you should be able to do now}
\begin{itemize}
  \item Implement a bit-serial CRC-15 engine matching CAN's polynomial.
  \item Explain why generating and checking are the same operation, and what a cleared CRC register
    at the end of a received frame proves.
  \item Work a few bits through the recursion by hand and compare against the simulation.
  \item Say why the engine has to be cleared between frames, and what failure arises if it is not.
\end{itemize}

\needspace{6\baselineskip}
\subsection*{Questions to test yourself with}
\begin{enumerate}
  \item Why does the engine need no mode for ``generate'' as against ``check''?
  \item What does it mean that the CRC register is zero after a received frame has been fed through
    it?
  \item What happens if a frame is aborted and the next frame starts being fed in without a
    \code{clear}?
  \item Why is a stuffed bit not fed into the CRC engine?
  \item Which signal in \code{can_controller} decides when \code{crc15} may update, and why is it not
    \code{bit_done}?
\end{enumerate}

\needspace{6\baselineskip}
\subsection*{Target}
By this point two of the four leaf modules should be finished. Which two does not matter.

\needspace{6\baselineskip}
\subsection*{Looking ahead}
\Chapref{c14:ch} builds the transmit shift register: serialising out onto the bus, with the same bit
stuffing rule as in \chapref{c10:ch}, now in VHDL. \Chapref{c14:ch} and \ref{c15:ch} together build
the modules deciding which bits reach \code{crc15} and when.
